\section{Minimum Weight Bi-Partition}
\label{sec:g-min-bipartition}
\problemlink{https://github.com/m-e-r-l-i-n/icpc-india/blob/main/src/ICPC\%202020\%20-\%20India\%20Prelims\%20Solutions.pdf}{Contest PDF (official 2020 solution slides; original CodeDrills page is offline)}
\source{ICPC 2020--21 Amritapuri Preliminary, problem 4 (PYQ)}{Medium}

\begin{reconbox}
Only the official solution slides survive. They reduce the problem to: ``a
tree is a connected bipartite graph, so after adding new edges it is enough to
find an MST; sorting vertices by $A_i$ and joining consecutive ones gives the
only $n-1$ new edges that matter; existing edges have weight $0$.'' This
section solves that core task. The full original statement may add conditions
that this reduction already handles; check the original if you find it.
\end{reconbox}

\problemstatement{A graph has $n$ vertices with values $A_i$ and $m$ existing
edges, which are free. You may add an edge $(u,v)$ at cost $|A_u-A_v|$. Find
the minimum total cost to make the graph connected (equivalently, the weight
of a minimum spanning tree of the complete graph where existing edges weigh
$0$).}

\begin{patternbox}
MST on a complete graph with weights $|A_u-A_v|$ $\Rightarrow$ only
\textbf{adjacent-in-sorted-order} edges can be useful: $O(n)$ candidate edges
instead of $O(n^2)$.
\end{patternbox}

\subsection*{Approach and intuition}

Sort the vertices by value: $A_{v_1}\le\dots\le A_{v_n}$. For $i<j$ the edge
$(v_i,v_j)$ has weight $A_{v_j}-A_{v_i}$, which equals the total weight of the
path $v_i,v_{i+1},\dots,v_j$, and each edge on that path is no heavier. By the
cycle property of MSTs, a heaviest edge on a cycle is never needed, so the
long edge can be dropped. The candidate set is: the $m$ free edges plus the
$n-1$ sorted-neighbour edges. Run Kruskal with a DSU.

\subsection*{Algorithm}
\begin{algorithmic}[1]
\State $E\gets\{(0,u,v)\text{ for existing edges}\}\cup\{(A_{v_{i+1}}-A_{v_i},v_i,v_{i+1})\}$
\State sort $E$; DSU on $n$ vertices; $cost\gets0$
\For{$(w,u,v)\in E$} \If{\textsc{Unite}$(u,v)$} $cost\gets cost+w$ \EndIf \EndFor
\State \Return $cost$
\end{algorithmic}

\dryrun{$A=[5,1,9,4]$, existing edge $(1,3)$. Sorted: $v=2,4,1,3$ (values
$1,4,5,9$). Edges: $(1,3)$ free; $(2,4)=3$, $(4,1)=1$, $(1,3)=4$. Kruskal:
free $(1,3)$, then $(4,1)$ cost $1$, then $(2,4)$ cost $3$; $(1,3)$ now
redundant. Total $4$. \checkmark\ (checked against Prim on the full graph)}

\subsection*{C++17 implementation}
\cppfile{greedy/28_min_weight_bipartition.cpp}

\begin{complexitybox}
\textbf{Time:} $O((n+m)\log(n+m))$. \textbf{Space:} $O(n+m)$.
\end{complexitybox}

\begin{pitfallbox}
\begin{itemize}
  \item Building all $\binom n2$ edges is $O(n^2)$ --- TLE for $n=10^5$.
  \item Self-loops and multi-edges among the free edges are harmless in
        Kruskal.
\end{itemize}
\end{pitfallbox}
